\section{A positivity toy: interaction signs and zero confinement}
\label{sec:exhibit-passivity}

A recurring motif in the framework is that the most rigidifying feasibility conditions are positivity conditions.
In physics, passivity constraints strongly restrict the analytic structure of response functions \cite{Willems1972II}; in mathematics, stability requirements often take the form of nonnegative quadratic forms, positive-definite kernels, or sign conditions on coefficients.
This exhibit gives a small, exactly understood instance: as a sign constraint on interactions is tightened, the zeros of a partition polynomial move onto a symmetry locus.

\subsection{The polynomial and its symmetries}
Take $n=\PassivityN$ spins $\sigma_i\in\{-1,1\}$ on a cycle with edge couplings $J_e$ and inverse temperature $\beta=\PassivityBeta$, and let
\[
  Z(z)\;=\;\sum_{\sigma\in\{-1,1\}^n}\exp\Bigl(\beta\sum_{e=\{i,j\}}J_e\,\sigma_i\sigma_j\Bigr)\,z^{\#\{i:\ \sigma_i=-1\}}\;=\;\sum_{k=0}^{n}a_k z^k,
\]
computed by direct enumeration.
Every coefficient $a_k$ is positive, whatever the signs of the $J_e$.
Reversing all spins leaves the interaction energy unchanged and exchanges $k$ with $n-k$, so $a_k=a_{n-k}$: the polynomial is palindromic as a matter of mathematics.
(The computation symmetrizes the coefficients numerically, which only removes floating-point summation asymmetry.)

A real palindromic polynomial has its roots closed under $z\mapsto 1/z$ and under complex conjugation, hence under the combined anti-involution $z\mapsto 1/\bar z$.
The fixed-point set of that anti-involution is the unit circle $|z|=1$, which is the natural symmetry locus.
(The inversion $z\mapsto 1/z$ alone fixes only $\pm 1$.)
Neither symmetry forces the roots onto the circle, and neither does positivity of the coefficients: $1+3z+z^2$ is real, palindromic, and positive, yet its two roots are real, negative, and off the circle.

\subsection{The feasibility parameter}
The couplings are drawn independently from a normal distribution with mean $0$ and standard deviation $\PassivityJScale$ (seed $0$), so they have mixed signs.
We tighten their signs with the clipped family
\[
  J_e(\lambda)=\max(J_e,0)+\lambda\min(J_e,0),\qquad\lambda\in[0,1],
\]
so that $\lambda=1$ is the original mixed-sign system and $\lambda=0$ removes every negative interaction.
This is an interaction-sign constraint; it is not a frequency-domain passivity condition, and we use ``positivity'' in that narrow sense.

At the endpoint the zeros are located by a classical theorem.
For $\beta\ge 0$ all effective couplings $\beta J_e(0)$ are nonnegative, and the Lee--Yang circle theorem \cite{LeeYang1952} applies to this finite ferromagnetic Ising polynomial, including zero couplings and disconnected components: every zero of $Z_0$ lies on $|z|=1$.
This is an imported theorem, not a numerical discovery.

A two-spin system shows the mechanism in closed form.
With one edge and $K=\beta J$, the polynomial is, up to a positive factor,
\[
  z^2+2e^{-2K}z+1 .
\]
If $K\ge 0$ the middle coefficient lies in $(0,2]$, the roots are complex conjugates with product $1$, and both lie on the unit circle.
If $K<0$ the middle coefficient exceeds $2$, and the roots are two negative reals with product $1$, one inside and one outside the circle.
Shrinking a negative coupling toward zero moves the coefficient toward $2$ and the roots toward $-1$.
This exact example supplies the motif that tightening the sign constraint \emph{can} confine zeros; no monotonicity theorem for arbitrary graphs is claimed.

\subsection{Measurement}
We compute all zeros of $Z_\lambda$ and measure their radial deviations $d_i=\bigl||z_i|-1\bigr|$, summarized by $\mathrm{mean\_dev}(\lambda)=\frac1n\sum_id_i$ and $\mathrm{max\_dev}(\lambda)=\max_id_i$.
The computation respects the structure of the problem.
Partition functions of disjoint subsystems multiply, so we factor $Z_\lambda$ over the connected components of the graph of nonzero couplings, with edges removed only when a coupling is exactly zero, and find the roots of each factor separately.
An isolated spin contributes exactly $1+z$.
Every component factor is palindromic, and a palindromic polynomial of odd degree has $-1$ as a root.
At $\lambda=0$ the components have degrees $\PassivityComponentsLamZero$, so $z=-1$ is a root of multiplicity four: one from each of the two cubic factors and one from each of the two isolated spins.
Finding the roots of the full degree-$12$ polynomial instead would split that fourfold root numerically and report spurious deviations of order $10^{-4}$; a small backward residual of the computed roots does not control the forward position of a repeated root.

\subsection{Results}
Table~\ref{tab:passivity} and Figure~\ref{fig:passivity-toy} show the outcome.
The mean deviation decreases at every sampled step, from $\PassivityMeanDevLamOne$ at $\lambda=1$ (where the largest deviation is $\PassivityMaxDevLamOne$) to $\PassivityMeanDevLamZero$ at $\lambda=0$, which is floating-point rounding: at the endpoint every zero is on the circle, as the Lee--Yang theorem requires.
As an independent check, the zeros at $\lambda=1$ and $\lambda=0$ were recomputed at $\PassivityCheckDPS$ digits with a separate enumeration and component decomposition.
The largest distance between matched roots is $\PassivityMatchLamOne$ at $\lambda=1$ and $\PassivityMatchLamZero$ at $\lambda=0$.
This high-precision computation is a consistency check, not an interval certificate; the endpoint statement rests on the theorem.

% Auto-generated by scripts/export_results_tex.py from notes/math_review/
\begin{table}[htbp]
\centering
\caption{Positivity toy: radial deviation of the zeros from the unit circle. The second column lists the degrees of the component factors of $Z_\lambda$; at $\lambda=0$ the deviations are at floating-point roundoff.}
\label{tab:passivity}
\small
\begin{tabular}{rlrr}
\toprule
$\lambda$ & Component degrees & mean\_dev & max\_dev \\
\midrule
$1$ & 12 & $0.872$ & $7.11$ \\
$0.7$ & 12 & $0.580$ & $4.27$ \\
$0.4$ & 12 & $0.350$ & $2.27$ \\
$0.2$ & 12 & $0.201$ & $1.22$ \\
$0.1$ & 12 & $0.0947$ & $0.718$ \\
$0$ & 3,3,4,1,1 & $1.55\times 10^{-15}$ & $6.00\times 10^{-15}$ \\
\bottomrule
\end{tabular}
\end{table}


\begin{figure}[htbp]
\centering
\includegraphics[width=\linewidth]{fig_passivity.pdf}
\caption{Positivity toy. (a) Zeros of $Z_\lambda$ plotted by argument and modulus, so that the unit circle is the horizontal line $|z|=1$; reciprocal pairs appear symmetric about it on the logarithmic scale. With mixed-sign couplings ($\lambda=1$) several zeros lie on the negative real axis (argument $\pi$) far from the circle; with nonnegative couplings ($\lambda=0$) all lie on it. (b) Mean and maximum radial deviation as the negative couplings are scaled to zero; the endpoint is at rounding level. Data in Table~\ref{tab:passivity}.}
\label{fig:passivity-toy}
\end{figure}

\subsection{Interpretation}
The endpoint is a theorem and the decrease across the five sampled steps in $\lambda$ is a finite observation for one sampled system.
The exhibit's role in this paper is to show a recognizable pattern under explicit control: a symmetry (\Ptwo) singles out a locus, a sign constraint acts as the feasibility condition, and the radial deviations of the zeros serve as the ledger (\Psix).
It suggests why problems of confining zeros to a self-dual locus invite positivity formulations, and it complements the prime exhibit, where staging and symmetrization without any additional feasibility condition did not yield a stable closure.
It establishes nothing about $\zeta$.
